
% =====================================================================
\section{Introduction}
\label{sec:intro}
% =====================================================================

\subsection{The question}

Fermions that attract one another pair, and the pairs condense. This is the origin of superconductivity
in metals \cite{BCS1957}, of the superfluid phases of liquid $^3$He \cite{Leggett1975}, and of the
superfluidity of nuclear matter. A condensate of pairs is described by a boson field, the expectation
value of a product of two fermion field operators. If the pair has no spin, the condensate is a scalar,
and scalar fields that fill space are what several open problems of physics call for: the Higgs field
of the Standard Model is one, and many proposals for the dark matter and for the discrepancy between
the measured values of the Hubble constant, the \emph{Hubble tension}, involve another.

An earlier study \cite{NashNCSS} asked whether the known neutrinos, in ordinary spacetime with three
space dimensions and one time, can form such a spin-singlet superfluid, and whether that state could
account for the dark matter or relieve the Hubble tension. The answer was negative on both counts. The
present paper and its companion \SIG{\cite{Companion1671}}{\cite{Companion1644}} take up the same
questions for a fermion field of a different kind: a field with sixteen components on a flat spacetime
of eight dimensions.

Eight dimensions are distinguished by an algebraic fact. The Dirac matrices of a spacetime are, by
definition, matrices $\gamma^A$ whose anticommutators reproduce the metric. In eight dimensions the
smallest such matrices have sixteen rows and columns, and for certain signatures of the metric they
can be chosen \emph{real}. Two such signatures are the subject of this work. The first has four
space-like and four time-like dimensions, the signature $4+4$; it is the case of the split octonion
algebra and of the triality of the group $SO(4,4)$, which the author has used before
\cite{Nash1990,Nash2010}. The second has seven space-like dimensions and one time-like dimension,
the signature $7+1$; it is the case of the octonions
\cite{Bryant1982,HarveyLawson1982,BrownGray1967}. (Real matrices of this size also exist for the
Euclidean signature and for five dimensions of one kind and three of the other
\cite{AtiyahBottShapiro1964,LawsonMichelsohn1989}.) A real set of Dirac matrices is the natural starting
point for a theory in which the sixteen components of the fermion field have a geometrical meaning,
and it is the starting point here.

%%<1644>
This paper treats the signature $4+4$. A spacetime with four times is not a small modification of
ordinary spacetime: its wave equation is not hyperbolic, and a field theory on it has no reason to be
stable \cite{Tegmark1997,CraigWeinstein2009,Bars2001}. Part of our task is therefore to find out which
part of the theory makes sense as a quantum theory, and to say so plainly. The companion paper
\cite{Companion1671} treats the signature $7+1$, which has one time and none of these difficulties.
The two papers follow the same plan, use the same methods and the same programs, and share the
sections that do not depend on the signature; where the two spacetimes lead to different conclusions,
the difference is stated.
%%</1644>
%%<1671>
This paper treats the signature $7+1$. It has one time, and its quantum field theory is conventional;
what is unconventional is that the four extra dimensions are space-like, and that the real Dirac
matrices force a particular sign convention on the metric and leave no real chirality. The companion
paper \cite{Companion1644} treats the signature $4+4$, in which four dimensions are time-like and the
theory makes sense as a quantum theory only in a restricted sector. The two papers follow the same
plan, use the same methods and the same programs, and share the sections that do not depend on the
signature; where the two spacetimes lead to different conclusions, the difference is stated.
%%</1671>

Three hypotheses are examined. They are stated here as hypotheses, and the paper reports what was found
for each, favorable or not.
\begin{itemize}
\item[H1.] A 16-component fermion field on this spacetime forms a superfluid state of spin 0 in which
the fermions are paired by the exchange of a new scalar boson, through a neutral current of the kind
known from the weak interaction.
\item[H2.] This superfluid state is the Higgs field.
\item[H3.] This superfluid state explains the observations of dark matter, or the Hubble tension, or
both.
\end{itemize}

\subsection{Objectives}
\label{sec:objectives}

The work has six objectives.
\begin{itemize}
\item[O1.] To fix the eight real $16\times16$ Dirac matrices of the spacetime without ambiguity, and to
prove their properties: that they are real, that they satisfy the anticommutation relations of the
metric, which group their scaled commutators generate, and what chirality means for them.
\item[O2.] To write down a valid Dirac equation for a 16-component field whose components are
complex-valued Grassmann operators, to determine in what sense it can be quantized, and to write it as
two coupled equations for 8-component fields by means of the chiral projections.
\item[O3.] To invent the scalar boson that mediates the pairing. We name it the \emph{Vorton-Scalar}.
The task is to define it by a Lagrangian, to construct its neutral current, and to derive its field
equations and its coupling to the fermions.
\item[O4.] To construct the superfluid state of spin 0: the pairing interaction, the gap equation, the
order parameter, the ground state, its thermodynamics and its collective modes; and to do this with the
effective action for composite operators of Cornwall, Jackiw and Tomboulis \cite{CJT1974}, which also
yields predictions for the unpaired state, for finite temperature and for the vacuum.
\item[O5.] To test hypothesis H2 against the conditions that a Higgs field must meet.
\item[O6.] To test hypothesis H3 by numerical simulation of the cosmological consequences of the
model.
\end{itemize}

\subsection{Strategy}

Four principles govern the work.

\emph{Exact algebra by machine.} Every matrix of this paper has the entries $0$, $+1$ and $-1$. An
identity between such matrices can be verified with integer arithmetic, and such a verification is a
proof. A program written for this project performs {{CHECKS}} verifications of this kind; each has an
identifier, and the statements of Secs.~\ref{sec:matrices} to \ref{sec:pairs} quote the identifiers on
which they rest. A second program, written independently in another language, repeats the essential
ones by different methods. Where an argument is short it is given in the text as well, so that the
reader can follow the logic without the computer.

\emph{One model, stated completely.} The Vorton-Scalar is defined by a Lagrangian with three
parameters. Everything else, including the mass of the fermions, the strength of the pairing force and
the cosmological behavior, is derived from it. Each approximation is named where it is made.

\emph{Two formulations of the many-body problem.} The paired state is first obtained from the gap
equation and then, independently, from the minimum of the effective potential of Cornwall, Jackiw and
Tomboulis, with different algorithms. The two agree to the accuracy of the arithmetic
(Sec.~\ref{sec:cjt}).

\emph{Simulation with validated tools.} The cosmological tests use an Einstein-Boltzmann solver that was
written for this project and validated against the established code \textsc{class}
\cite{BlasLesgourguesTram2011} (Sec.~\ref{sec:hubble}).

All programs, all data and the documents that record the provenance of every number are public
\cite{SpacetimesRepo}. No computed result in this paper was copied by hand: the text is assembled by
a program that inserts each such number from the data file in which the computation stored it. Values
quoted from the literature and the parameter values of the examples are written in the text.

\subsection{Summary of results}

%%<1644>
\emph{The matrices and the field (O1, O2).} The eight Dirac matrices are signed permutation matrices
built from the multiplication table of the split octonions. Their scaled commutators span the Lie
algebra $\so(4,4)$ and, together with the matrices themselves, generate the group $\Pin(4,4)$, which
has four connected components. The chirality matrix is real, and the chiral projectors select the
upper and the lower eight components. The Dirac equation contains no factor $i$. Its evolution in the
ordinary time is unitary, and a positive conserved density exists, if and only if the field does not
depend on the three additional times; fields that do depend on them include exponentially growing
modes. We define the quantized theory by that restriction and state it as an assumption wherever it is
used.

\emph{The pair (O4).} A single 16-component field has no Lorentz-invariant pair. It has a vector of
pairs, and the component of that vector along the hidden space-like coordinate $x0$ is a scalar under
the Lorentz group of ordinary spacetime. It is the only $s$-wave pair that is invariant both under
ordinary rotations and under the rotations of the three additional times.

\emph{The Vorton-Scalar (O3).} In the fully covariant theory on $4+4$ spacetime the construction is
impossible: the chiral phase rotation is a symmetry of the mass term and not of the kinetic term, the
reverse of the familiar situation. In the quantized sector it exists. The Vorton-Scalar is a complex
scalar field whose vacuum value gives the fermions their mass. Its phase mode is massless, and its
exchange is a current-current interaction between the neutral current of the scalar and the axial
current of the fermions, with a Fermi constant fixed by the vacuum value.

\emph{The superfluid (O4).} The radial mode of the Vorton-Scalar attracts two fermions in every pair
channel; the neutral current attracts in the four vector channels along the extra coordinates and
repels in the others. The gap of the vector channel is the largest. In the sector of fields that
depend on ordinary space only, the order parameter is a four-vector of a hidden $SO(4)$ symmetry, which
is not a symmetry of the spacetime. The hidden-axis singlet is one of a family of degenerate ground
states; the degeneracy is exact in that sector and is lifted only by a dependence of the fields on
$x0$. The effective action of Cornwall, Jackiw and Tomboulis reproduces the gap, the critical
temperature, the critical coupling of the two-body bound state and the Ginzburg-Landau coefficients,
and it predicts that the fermion mass vanishes in dense matter, that matter is self-bound for strong
coupling, that the mass melts at a temperature of the order of the vacuum value, and that a condensate
in empty space requires a strong coupling.

\emph{The Higgs hypothesis (O5).} The pair is a scalar of ordinary spacetime. The breaking pattern of
the Higgs vacuum, $SO(4)\to SO(3)$, is present only as the hidden symmetry, which the geometry cannot
turn into a gauge symmetry; under the rotations of the additional times, the only non-Abelian internal
symmetry that the geometry of $4+4$ provides, the condensate is a singlet and cannot break the
electroweak symmetry. The hypothesis has no geometric basis in this spacetime.
%%</1644>
%%<1671>
\emph{The matrices and the field (O1, O2).} The eight Dirac matrices are signed permutation matrices
built from the multiplication table of the octonions. Real matrices exist only if the time-like matrix
squares to $+1$ and the seven space-like ones to $-1$. Their scaled commutators span the Lie algebra
$\so(7,1)$ and, together with the matrices themselves, generate the group $\Pin(7,1)$, which has four
connected components. No real chirality matrix exists; the chiral projectors are complex. The Dirac
equation necessarily contains a factor $i$ in front of its real matrices. Its Hamiltonian is real and
symmetric, the evolution is unitary, and the field is quantized in the conventional way. A particle
has eight internal states.

\emph{The pair (O4).} A single 16-component field has no Lorentz-invariant pair, and no $s$-wave pair
that is invariant under all rotations of the seven space dimensions. It has a vector of pairs, and the
component of that vector along the hidden space-like coordinate $x0$ is a scalar under the Lorentz
group of ordinary spacetime. It is the only $s$-wave pair that is invariant both under ordinary
rotations and under the rotations of the three other extra coordinates.

\emph{The Vorton-Scalar (O3).} It is a complex scalar field whose vacuum value gives the fermions
their mass, so that the chiral rotation of the fermions, accompanied by a phase rotation of the scalar,
is an exact symmetry. Its phase mode is massless, and its exchange is a current-current interaction
between the neutral current of the scalar and the axial current of the fermions, with a Fermi constant
fixed by the vacuum value.

\emph{The superfluid (O4).} The radial mode of the Vorton-Scalar attracts two fermions in every pair
channel; the neutral current attracts in the four vector channels along the extra coordinates and
repels in the others. The gap of the vector channel is the largest. The order parameter is a
four-vector of the rotation group $SO(4)$ of the four extra space dimensions, and its direction is
chosen spontaneously: the hidden-axis singlet is one of a family of degenerate ground states. The
effective action of Cornwall, Jackiw and Tomboulis reproduces the gap, the critical temperature, the
critical coupling of the two-body bound state and the Ginzburg-Landau coefficients, and it predicts
that the fermion mass vanishes in dense matter, that matter is self-bound for strong coupling, that the
mass melts at a temperature of the order of the vacuum value, and that a condensate in empty space
requires a strong coupling.

\emph{The Higgs hypothesis (O5).} The pair is a scalar of ordinary spacetime, it belongs to the
representation $(2,2)$ of $SU(2)\times SU(2)$, and its condensate breaks $SO(4)$ to $SO(3)$: the
representation and the breaking pattern of the Higgs field and of its custodial symmetry. This is a
necessary condition, and it is met. None of the dynamical conditions has been shown, and the simplest
estimates of two of them are unfavorable.
%%</1671>

\emph{Cosmology (O6).} The model contains dark radiation of exactly the kind that has been proposed as
a remedy of the Hubble tension, with a step in its density when the fermions annihilate. For
microwave-background data of the quality of the Planck satellite that agree with the standard model,
raising the Hubble constant to the locally measured value costs at least {{GEN:dr_sig_local_min_all}}
standard deviations, and {{GEN:dr_sig_local_min_model}} or more if the step of the model falls into the
epochs that the data probe. Treated as dark matter, the superfluid suppresses the growth of structure;
the fermion masses that structure formation allows are too large for the fermions to be degenerate in
any galaxy, and a condensate of bound pairs cannot support cores larger than {{GEN:sfdm_rtf_pc_5p3}}
parsec. The superfluid explains neither the Hubble tension nor the galactic phenomenology of dark
matter. Heavy bound pairs or clusters of the fermions remain a possible, and so far untestable, form of
cold dark matter.

\subsection{Organization, units and notation}
\label{sec:notation}

Section~\ref{sec:matrices} defines the Dirac matrices and proves their properties.
Section~\ref{sec:dirac} gives the Dirac equation and its quantization. Section~\ref{sec:pairs}
determines which fermion pairs exist. Section~\ref{sec:vs} invents the Vorton-Scalar.
Section~\ref{sec:pairing} derives the pairing interaction, Sec.~\ref{sec:cjt} the paired state, and
Sec.~\ref{sec:order} its order parameter and ground state. Section~\ref{sec:matter} follows the
Vorton-Scalar into matter, to finite temperature and into the vacuum. Section~\ref{sec:higgs} examines
the Higgs hypothesis, Secs.~\ref{sec:hubble} and \ref{sec:dm} the Hubble tension and the dark matter.
Section~\ref{sec:predictions} collects the predictions, Sec.~\ref{sec:status} lists the principal
statements with their status, and Sec.~\ref{sec:concl} concludes. Each of
Secs.~\ref{sec:matrices} to \ref{sec:dm} ends with a subsection that separates what has been
established from what has not.

\emph{Coordinates.} The eight coordinates are
\begin{equation}
X=\{x0,\;x1,\;x2,\;x3,\;x4,\;x5,\;x6,\;x7\}.
\label{eq:coords}
\end{equation}
They are never renamed. The coordinates $x1,x2,x3$ are the Cartesian coordinates $x,y,z$ of ordinary
space, and $x4$ is the ordinary time. The coordinate $x0$ is a \emph{hidden} space-like coordinate.
The coordinates $x5,x6,x7$ are \SIG{three additional time-like coordinates}{three additional
space-like coordinates}. A capital index $A,B,\dots\in\{0,\dots,7\}$ labels the coordinate $xA$, with
$\partial_A=\partial/\partial x^A$ and summation over repeated indices. Three smaller index sets recur:
$i,j\in\{1,2,3\}$ for ordinary space; $a,b\in\{0,5,6,7\}$ for the hidden axis and the three additional
coordinates, which we call \emph{the four extra axes}; and $t\in\{5,6,7\}$ for the three additional
coordinates alone.

\emph{Units.} $\hbar=c=k_B=1$.

\emph{Tags.} The companion paper uses the same letters for the corresponding objects of the other
spacetime, which are different matrices. To keep the two apart, every object of the $4+4$ spacetime
carries the tag 1644 and every object of the $7+1$ spacetime the tag 1671 (16 components; signature
$4+4$ or $7+1$): $\gamma_{1644}^A$ and $\gamma_{1671}^A$, for example. In this paper the tag
$\tg$ is written in the equations that define an object or state a result, and omitted in derivations;
an object of the other spacetime always carries its tag. The 16-component fields of the two papers are
named \texttt{dirac1644} and \texttt{dirac1671}.

\emph{Two stages.} The spacetime of this paper is flat. In the program to which the paper belongs
this is the first of two stages. In the second stage, which is not implemented here, the metric is
dynamical: the three ordinary space dimensions inflate, the three additional dimensions deflate, which
is the proposed reason why they have not been detected, and the hidden coordinate $x0$ curls up into a
ring of very small radius. Where a result of this paper depends on what the second stage would have
to provide, this is said.
